\section{Evidence from public data}\label{sec:evidence}

\subsection{USDT on Tron is a retail dollar rail}\label{sec:rails}

\Cref{fig:supply} shows the migration of USDT onto Tron. At the end of 2019 Tron carried 19\% of the USDT supply on the
covered chains; by mid-2021 it carried half, and on 23 May 2026 it carried 47\% (US\$89 billion) against 51\% on
Ethereum (US\$97 billion). The share has oscillated around one half since 2021, so ``Tron versus Ethereum'' is not a
story of one chain displacing the other but of two rails of similar size.

\begin{figure}[!htbp]
\centering
\fig{0.98}{fig_rails_supply}
\caption{USDT supply by blockchain. Top: US\$ billions, stacked. Bottom: Tron's share of the covered supply.}
\label{fig:supply}
\end{figure}

The two rails serve different clienteles. \Cref{tab:onchain} reports a snapshot and \cref{fig:tworails} the history. On
23 May 2026, 74.5 million addresses held USDT on Tron, 5.3 times the 14.0 million on Ethereum, and the mean balance per
holding address was about US\$1,200 against about US\$6,900. Tron recorded 2.1 million USDT transfers that day, 3.1
times Ethereum's 0.67 million, with less supply. Measured as transfers per US\$1 million of supply---a count-based
turnover rate---Tron's median since January 2024 is 32 a day, 9.3 times that of USDT on Ethereum (3.4) and 6.7 times
that of USDC on Ethereum (4.8). Small balances turned over often are the signature of a payments rail; large balances
turned over rarely are the signature of a settlement and treasury rail.

\begin{table}[!htbp]
\centering\small
\begin{threeparttable}
\caption{On-chain snapshot of the main dollar-stablecoin rails, 23 May 2026}
\label{tab:onchain}
\begin{tabular}{lrrrrrrr}
\toprule
 & Supply & 1-yr & Holders & 1-yr & Active & Transfers & Mean \\
 & US\$bn & \% & m & \% & k/day & k/day & balance US\$ \\
\midrule
USDT on Tron & 89.36 & +21 & 74.49 & +13 & 941 & 2,071 & 1,200 \\
USDT on Ethereum & 97.07 & +33 & 14.01 & +83 & 212 & 667 & 6,926 \\
USDC on Ethereum & 52.72 & +31 & 6.99 & +118 & 118 & 505 & 7,543 \\
USDC on Tron & 0.03 & -30 & 0.38 & +0 & 0 & 0 & 74 \\
\bottomrule
\end{tabular}
\begin{tablenotes}[flushleft]\footnotesize
\item \emph{Notes:} Supply in US\$ billions; holders are addresses with a positive balance (millions); active addresses
and transfers are daily counts (thousands); 1-yr columns give percentage changes over 365 days. Mean balance is supply
divided by holders. Source: Coin Metrics Community Data; author's calculations.
\end{tablenotes}
\end{threeparttable}
\end{table}

\begin{figure}[!htbp]
\centering
\fig{0.98}{fig_two_rails}
\caption{Holders, mean balances and turnover of USDT on Tron and Ethereum and of USDC on Ethereum. 30-day rolling
medians; the middle and right panels use logarithmic scales.}
\label{fig:tworails}
\end{figure}

Two caveats apply. An address is not a person: exchanges hold customer funds in omnibus wallets, and individuals hold
many addresses. Both biases are, if anything, larger on Tron, where exchanges custody much of the retail float and
address creation is cheap, so the comparison of balances per address is indicative rather than exact. The turnover
comparison does not depend on address counts and is the more robust of the two.

\subsection{Circle's exit left the cheapest rail to a single issuer}\label{sec:usdcexit}

USDC on Tron grew to US\$3.5 billion in August 2022 and then shrank. After Circle announced that it would stop minting
USDC on Tron in early 2024, supply fell below US\$30 million (\cref{fig:usdcexit}). For users in Southeast Asia who
need a cheap, liquid dollar token, USDT on Tron is now, in practice, the only option at scale. This concentration cuts
both ways for enforcement: one issuer with the technical ability to freeze tokens sits on the main rail, but
regulatory regimes that would restrict USDT, such as the European Union's, have no regional substitute to push users
towards.

\begin{figure}[!htbp]
\centering
\fig{0.98}{fig_usdc_exit}
\caption{USDC and USDT supply on the Tron blockchain, logarithmic scale. The shaded band marks the approximate
wind-down period.}
\label{fig:usdcexit}
\end{figure}

\subsection{Thailand has little macroeconomic reason to hold dollars on-chain}\label{sec:macro}

Where inflation is high, currencies are volatile or capital is controlled, people hold dollar stablecoins for reasons
that have nothing to do with crime. \Cref{fig:macro} shows that Thailand is not such a place. Its consumer-price
inflation averaged 1.7\% a year over 2019--2023, the fourth lowest of 23 economies after Japan, China and Malaysia. The
baht's volatility against the dollar since January 2021, 7.8\% a year, is in the middle of the 21 currencies in the
Federal Reserve's H.10 release---below the yen, the won and the Australian dollar.

The neighbourhood is different. Inflation averaged 14.9\% in Laos over 2019--2024 and was 8.8\% in Myanmar in 2019, the
last year the World Bank reports, before a period of currency collapse and capital controls. Cambodia is largely
dollarised already. Legitimate dollar demand on the other side of Thailand's borders can and plausibly does flow
through Thai platforms and peer-to-peer markets, especially given the millions of migrant workers from these countries
in Thailand. This is the most important legitimate confounder for any claim that Thai stablecoin flows are
``excessive'', and the pre-registered plan treats it as such (H1 and H6 in \cref{tab:prereg}).

\begin{figure}[!htbp]
\centering
\fig{0.98}{fig_macro}
\caption{(a) Consumer-price inflation, mean of available years 2019--2024 (dot) and range (bar); hollow markers denote
fewer than three observations. (b) Annualised volatility of daily log changes in each currency's exchange rate against
the US dollar, January 2021 to October 2026.}
\label{fig:macro}
\end{figure}

\subsection{Sanctions lists point to Tron}\label{sec:ofac}

\Cref{tab:ofac} classifies the 1,073 digital-currency addresses on the OFAC list by the currency tag OFAC assigns and by
address encoding. Of the 368 addresses tagged USDT, TRX or USDC, 351 (95\%) are Tron addresses; among the 94 tagged
USDT, 79 are Tron addresses, eight are Ethereum-style addresses and seven are Bitcoin-style addresses that operated on
the legacy Omni layer. Counts of listed addresses are a crude measure: one listing can cover a single wallet or the hub
of a large network, and listings reflect enforcement priorities. They nonetheless confirm that, for the authority with
the most active crypto-sanctions programme, the stablecoin activity worth sanctioning happens mainly on Tron.

\begin{table}[!htbp]
\centering\small
\begin{threeparttable}
\caption{OFAC-listed digital-currency addresses by currency tag and address encoding}
\label{tab:ofac}
\begin{tabular}{lrrrrr}
\toprule
Listed under & Bitcoin & EVM (0x\ldots{}) & Other & Tron (T\ldots{}) & Total \\
\midrule
XBT & 531 & 0 & 0 & 1 & 532 \\
TRX & 0 & 0 & 0 & 272 & 272 \\
ETH & 0 & 123 & 0 & 0 & 123 \\
USDT & 7 & 8 & 0 & 79 & 94 \\
LTC & 0 & 0 & 13 & 0 & 13 \\
XMR & 0 & 0 & 11 & 0 & 11 \\
BCH & 1 & 0 & 6 & 0 & 7 \\
DASH & 0 & 0 & 5 & 0 & 5 \\
\bottomrule
\end{tabular}
\begin{tablenotes}[flushleft]\footnotesize
\item \emph{Notes:} The eight currency tags with the most listed addresses; the remaining nine tags hold 34 addresses.
Encoding is inferred from address syntax. Source: OFAC SDN list via the \texttt{0xB10C} mirror, updated
7 November 2026; author's classification.
\end{tablenotes}
\end{threeparttable}
\end{table}

\subsection{What chain-wide data cannot do}\label{sec:eventstudy}

It is tempting to evaluate crackdowns with the free data used above. For each of three shocks dated to the day---the
power cut to border scam hubs on 5 February 2025, the block on unlicensed offshore exchanges on 28 June 2025 and the
Prince Group sanctions on 14 October 2025---I estimate a counterfactual for the log of daily Tron USDT transfers from a
linear trend with day-of-week effects fitted over days $-150$ to $-31$, and cumulate abnormal activity over days 0 to
30. Bands come from a moving-block bootstrap (14-day blocks) that redraws both the estimation residuals, so that the
uncertainty of the extrapolated trend is included, and the future shocks.

\begin{figure}[!htbp]
\centering
\fig{0.98}{fig_event_study}
\caption{Top: cumulative abnormal log transfers of USDT on Tron after three enforcement shocks, with moving-block
bootstrap 95\% bands; grey dots are daily abnormal values before the event. Bottom: distribution of the 30-day
cumulative abnormal value at 149 placebo dates at least 45 days from any of the three shocks.}
\label{fig:eventstudy}
\end{figure}

\Cref{fig:eventstudy} reports the result. Two of the three shocks appear ``significant'' with opposite signs: after the
offshore block, activity falls below the band (cumulative abnormal value $-2.83$ log points by day 30); after the
sanctions it rises above it ($+2.54$). The bottom panel shows why neither should be believed. Repeating the exercise at
149 placebo dates spaced a week apart, the nominal bands exclude zero at 58\% of them. Chain-wide activity is driven by
global market events---the sanctions date, for instance, falls days after one of the largest liquidation episodes in
crypto history---and by structural breaks that no trend model anticipates. Calibrated against the placebo distribution,
the shocks have $p$-values of 1.00, 0.69 and 0.73. The 95th percentile of the absolute placebo values is 11.1 log
points, so a shock would have to move global Tron USDT activity by about 36 log points a day, sustained for a month,
to stand out. A national crackdown cannot plausibly do that, and by \cref{prop:displacement} it need not change the
total at all.

The lesson is general. Evaluations of anti-laundering policy that rely on chain-wide aggregates, or on any series
dominated by global activity, should report placebo-calibrated inference, and null results from such data say little
about whether a policy worked. Credible evaluation needs data on the composition of flows through the channel that the
policy targets.
